\clearpage
\fancyhead[L]{\small GERLING. SOLAR ECLIPSE.}
\hypertarget{native-627-lower}{}
\begin{center}\small Göttingische gelehrte Anzeigen. No. 103. Pages 1025..1030. 26 June 1820.\\[3pt]\rule{.72\linewidth}{.3pt}\end{center}

Göttingen. Concerning the great solar eclipse that will occur on 7 September of this year, Professor \textsc{Gerling} in Marburg, who had already discussed it eight years ago in his inaugural dissertation (see our Anzeigen, 1812, p. 1889 [p. 542 of this volume]), has, in a letter to Court Councillor \textsc{Gauss} dated 23 May, communicated the results of his repeated and extended investigations. Since this eclipse will appear annular in most of Germany and will long be the most remarkable phenomenon of its kind in these regions, we believe we can do many readers of these pages \SourceLeaf{628} a welcome service by publishing an excerpt from those communications.

The following positions of the Sun and Moon, taken from \textsc{Carlini}'s solar tables and \textsc{Bürg}'s lunar tables, served as the basis for the calculation:

7 September 1820. Paris mean time.

\begin{center}
\small\setlength{\tabcolsep}{2pt}\begin{tabular}{l|rrr|rrr|rrr}
 & \multicolumn{3}{c|}{\shortstack{Longitude of\\the Sun.}} & \multicolumn{3}{c|}{\shortstack{Longitude of\\the Moon.}} & \multicolumn{3}{c}{\shortstack{Northern latitude\\of the Moon.}} \\
\hline
$11^{\mathrm u}$ a.m. & $164^\circ$ & $40'$ & $26''4$ & $163^\circ$ & $19'$ & $53''6$ & $0^\circ$ & $52'$ & $45''0$ \\
2 p.m. & & $47$ & $44.0$ & $164$ & $48$ & $17.7$ & & $44$ & $41.6$ \\
5 & & $55$ & $1.7$ & $166$ & $16$ & $38.3$ & & $36$ & $34.9$ \\
\end{tabular}
\end{center}

\begin{center}\small\setlength{\tabcolsep}{3pt}
\begin{tabular}{lr}
Horizontal parallax of the Moon . . . . . . . . . . & $53'\ 55''5$ \\
\hspace{1em}---\hspace{1em}of the Sun . . . . . . . . . . & $8.7$ \\
Semidiameter of the Moon . . . . . . . . . . . . . & $14'\ 43.1$ \\
\hspace{1em}---\hspace{1em}of the Sun . . . . . . . . . . . . . & $15\ 54.8$ \\
Obliquity of the ecliptic . . . . . . . . . . . . . & $23^\circ\ 27\ 56.2$ \\
\end{tabular}
\begin{tabular}{ll}
The conjunction therefore occurs at . . . . . . & $1^{\mathrm u}\ 58^{\mathrm m}\ 46^{\mathrm s}3$ mean time in Paris \\
at longitude . . . . . . . . . . . . . . . & $164^\circ\ 47'\ 41''$ \\
\end{tabular}
\end{center}

For the following individual locations, Professor \textsc{Gerling} has calculated the principal times of the phenomenon in advance to facilitate observations:

\begin{center}\small\setlength{\tabcolsep}{3pt}
\begin{tabular}{l|rrr|rrr|rrr}
 &  \multicolumn{3}{c|}{Beginning.}  &  \multicolumn{3}{c|}{Middle.}  &  \multicolumn{3}{c|}{End.}  \\
\hline
Berlin  &  $1^{\mathrm u}$  &  $31^{\mathrm m}$  &  $55^{\mathrm s}$  &  $2^{\mathrm u}$  &  $57^{\mathrm m}$  &  $31^{\mathrm s}$  &  $4^{\mathrm u}$  &  $15^{\mathrm m}$  &  $39^{\mathrm s}$  \\
Bremen  &  $1$  &  $7$  &  $13$  &  $2$  &  $33$  &  $38$  &  $3$  &  $54$  &  $5$  \\
Breslau  &  $1$  &  $53$  &  $36$  &  $3$  &  $18$  &  $29$  &  $4$  &  $34$  &  $54$  \\
Cassel  &  $1$  &  $13$  &  $39$  &  $2$  &  $40$  &  $30$  &  $4$  &  $0$  &  $54$  \\
Cöln  &  $1$  &  $1$  &  $6$  &  $2$  &  $28$  &  $52$  &  $3$  &  $50$  &  $16$  \\
Danzig  &  $1$  &  $55$  &  $32$  &  $3$  &  $19$  &  $5$  &  $4$  &  $34$  &  $22$  \\
Göttingen  &  $1$  &  $15$  &  $51$  &  $2$  &  $42$  &  $28$  &  $4$  &  $2$  &  $39$  \\
Gotha  &  $1$  &  $20$  &  $59$  &  $2$  &  $47$  &  $34$  &  $4$  &  $7$  &  $30$  \\
Halle  &  $1$  &  $26$  &  $26$  &  $2$  &  $52$  &  $20$  &  $4$  &  $11$  &  $45$  \\
Hamburg  &  $1$  &  $12$  &  $26$  &  $2$  &  $38$  &  $19$  &  $3$  &  $58$  &  $14$  \\
\end{tabular}
\end{center}
\begin{center}\small\setlength{\tabcolsep}{3pt}
\begin{tabular}{l|rrr|rrr|rrr}
  &  \multicolumn{3}{c|}{Beginning.}  &  \multicolumn{3}{c|}{Middle.}  &  \multicolumn{3}{c}{End.}  \\
\hline
 Hannover  &  $1^{\mathrm u}$  &  $13^{\mathrm m}$  &  $11^{\mathrm s}$  &  $2^{\mathrm u}$  &  $39^{\mathrm m}$  &  $36^{\mathrm s}$  &  $3^{\mathrm u}$  &  $59^{\mathrm m}$  &  $45^{\mathrm s}$ \\
 Königsberg  &  $2$  &  $4$  &  $43$  &  $3$  &  $27$  &  $20$  &  $4$  &  $41$  &  $32$ \\
 Leipzig  &  $1$  &  $28$  &  $38$  &  $2$  &  $54$  &  $25$  &  $4$  &  $13$  &  $45$ \\
 Mannheim  &  $1$  &  $11$  &  $47$  &  $2$  &  $39$  &  $36$  &  $4$  &  $0$  &  $31$ \\
 Marburg  &  $1$  &  $11$  &  $23$  &  $2$  &  $38$  &  $40$  &  $3$  &  $59$  &  $18$ \\
 München  &  $1$  &  $30$  &  $40$  &  $2$  &  $57$  &  $47$  &  $4$  &  $17$  &  $33$ \\
 Ofen  &  $2$  &  $11$  &  $2$  &  $3$  &  $35$  &  $49$  &  $4$  &  $51$  &  $31$ \\
 Prag  &  $1$  &  $41$  &  $51$  &  $3$  &  $7$  &  $50$  &  $4$  &  $25$  &  $46$ \\
 Tübingen  &  $1$  &  $16$  &  $42$  &  $2$  &  $44$  &  $36$  &  $4$  &  $5$  &  $24$ \\
 Wien  &  $1$  &  $55$  &  $45$  &  $3$  &  $21$  &  $24$  &  $4$  &  $38$  &  $28$ \\
\end{tabular}
\end{center}

The times here and in the following table are given in \emph{apparent solar time}.

\begin{center}\small\setlength{\tabcolsep}{3pt}
\begin{tabular}{l|r|r|rrr|rrr}
 &  $P$  &  $D$  &  \multicolumn{3}{c|}{\shortstack{First internal\\contact.}}  &  \multicolumn{3}{c|}{\shortstack{Second internal\\contact.}}  \\
\hline
Berlin  &  $72^\circ$  &  $84''$  &  ---  &  ---  &   &  ---  &  ---  &   \\
Bremen  &  $65$  &  $9$  &  $2^{\mathrm u}$  &  $30'$  &  $46''$  &  $2^{\mathrm u}$  &  $36^{\mathrm m}$  &  $30^{\mathrm s}$  \\
Breslau  &  $78$  &  $131$  &  ---  &  ---  &   &  ---  &  ---  &   \\
Cassel  &  $66$  &  $1$  &  $2$  &  $37$  &  $37$  &  $2$  &  $43$  &  $24$  \\
Cöln  &  $62$  &  $55$  &  $2$  &  $27$  &  $23$  &  $2$  &  $30$  &  $21$  \\
Danzig  &  $79$  &  $186$  &  ---  &  ---  &   &  ---  &  ---  &   \\
Göttingen  &  $67$  &  $10$  &  $2$  &  $39$  &  $37$  &  $2$  &  $45$  &  $19$  \\
Gotha  &  $68$  &  $18$  &  $2$  &  $44$  &  $47$  &  $2$  &  $50$  &  $20$  \\
Halle  &  $70$  &  $48$  &  $2$  &  $50$  &  $20$  &  $2$  &  $54$  &  $19$  \\
Hamburg  &  $66$  &  $36$  &  $2$  &  $35$  &  $55$  &  $2$  &  $40$  &  $43$  \\
\end{tabular}
\end{center}
\begin{center}\small\setlength{\tabcolsep}{3pt}
\begin{tabular}{l|r|r|rrr|rrr}
  &  $P$  &  $D$  &  \multicolumn{3}{c|}{\shortstack{First internal\\contact.}}  &  \multicolumn{3}{c}{\shortstack{Second internal\\contact.}}  \\
\hline
 Hannover  &  $66^\circ$  &  $17''$  &  $2^{\mathrm u}$  &  $36^{\mathrm m}$  &  $49^{\mathrm s}$  &  $2^{\mathrm u}$  &  $42^{\mathrm m}$  &  $23^{\mathrm s}$ \\
 Königsberg  &  $81$  &  $217$  &  ---  &  ---  &   &  & --- &  \\
 Leipzig  &  $71$  &  $52$  &  $2$  &  $52$  &  $49$  &  $2$  &  $56$  &  $1$ \\
 Mannheim  &  $64$  &  $44$  &  $2$  &  $37$  &  $28$  &  $2$  &  $41$  &  $43$ \\
 Marburg  &  $65$  &  $18$  &  $2$  &  $35$  &  $53$  &  $2$  &  $41$  &  $27$ \\
 München  &  $71$  &  $2$  &  $2$  &  $54$  &  $54$  &  $3$  &  $0$  &  $40$ \\
 Ofen  &  $84$  &  $123$  &  ---  &  ---  &   &  ---  &  ---  &  \\
 Prag  &  $75$  &  $75$  &  ---  &  ---  &   &  ---  &  ---  &  \\
 Tübingen  &  $66$  &  $46$  &  $2$  &  $42$  &  $35$  &  $2$  &  $46$  &  $37$ \\
 Wien  &  $79$  &  $87$  &  ---  &  ---  &   &  ---  &  ---  &  \\
\end{tabular}
\end{center}

Here $P$ denotes the point on the solar limb where the Moon first encroaches at the beginning of the eclipse, counting degrees from the northernmost point (with respect to the vertical circle) \SourceLeaf{629} westward. $D$, by contrast, gives in seconds of arc the apparent distance between the centers of the Sun and the Moon at mid-eclipse. In all calculations the flattening is taken to be $\tfrac{1}{300}$.

The principal times of the eclipse for other places in Germany can be derived approximately from those given in the first table above for a nearby place by means of the following formulas. In these formulas, d$L$ denotes the difference in longitude between the two places in seconds of time, d$B$ the difference in latitude in seconds of arc, and d$Z$ the adjustment to be added to the values in the table. Take d$L$ to be positive if the place for which the times are sought lies east of the place chosen from the table, and d$B$ likewise to be positive if the former lies north of the latter.

For the beginning
\[
\mathrm{d}Z = +1.313\,\mathrm{d}L - 0.0296\,\mathrm{d}B
\]
For mid-eclipse
\[
\mathrm{d}Z = +1.242\,\mathrm{d}L - 0.0343\,\mathrm{d}B
\]
For the end
\[
\mathrm{d}Z = +1.118\,\mathrm{d}L - 0.0363\,\mathrm{d}B
\]

It is especially important to know the boundary lines of the zone within which the eclipse appears annular. Prof. \textsc{Gerling} has determined by calculation the following points on these boundary lines, with longitudes reckoned from Ferro.

\begin{center}
\begin{tabular}{c|c|c|c|c}
\shortstack{Apparent\\time in\\Paris.} & \multicolumn{2}{c|}{\shortstack{Contact of the\\southern limbs}} & \multicolumn{2}{c}{\shortstack{Contact of the\\northern limbs.}} \\
 & Latitude. & Longitude. & Latitude. & Longitude. \\
\hline
$2^{\mathrm u}\ 4^{\mathrm m}$ & & & $52^\circ\ 59'$ & $22^\circ\ 43'$ \\
6 & & & $52\ \ 15$ & $23\ \ 15$ \\
8 & $54^\circ\ 12'$ & $28^\circ\ 50'$ & $51\ \ 31$ & $23\ \ 46$ \\
10 & $53\ \ 26$ & $29\ \ 18$ & $50\ \ 47$ & $24\ \ 17$ \\
12 & $52\ \ 40$ & $29\ \ 47$ & $50\ \ \ 4$ & $24\ \ 48$ \\
14 & $51\ \ 55$ & $30\ \ 16$ & $49\ \ 21$ & $25\ \ 19$ \\
16 & $51\ \ 10$ & $30\ \ 46$ & $48\ \ 39$ & $25\ \ 49$ \\
18 & $50\ \ 26$ & $31\ \ 15$ & $47\ \ 57$ & $26\ \ 19$ \\
20 & $49\ \ 42$ & $31\ \ 45$ & $47\ \ 15$ & $26\ \ 49$ \\
22 & $48\ \ 58$ & $32\ \ 15$ & $46\ \ 34$ & $27\ \ 20$ \\
24 & $48\ \ 15$ & $32\ \ 46$ & $45\ \ 53$ & $27\ \ 50$ \\
26 & $47\ \ 32$ & $33\ \ 17$ & & \\
28 & $46\ \ 50$ & $33\ \ 49$ & & \\
\end{tabular}
\end{center}

If these points are joined by straight lines on a map, one encounters the following places, none of which is more than about $2\tfrac{1}{2}$ Meilen from the boundary lines. The places marked by a change of type lie closest to the boundary.

For contact of the \emph{southern} limbs: Heiligenhafen, Neustadt, Grevesmühlen, Wismar, Schwerin, Parchim, Grabow, \emph{Perleberg}, Werben, Havelberg, \emph{Arneburg}, Tangermünde, Rathenow, Genthin, Ziesar, Coswig, Wittenberg, \emph{Kemberg}, Düben, Torgau, Wurzen, Oschatz, Döbeln, \emph{Freiberg}, Seyda, Kommotau, \emph{Rakonitz}, Przibram, \emph{Teyn}, Budweis, \emph{Grazen}, \emph{Ips}, \emph{Oberndorf}, Tirnitz, Fürstenfeld etc.

For contact of the \emph{northern} limbs: Staveren, Medenblick, Harderwyk, Amersfort, \emph{Barneveld}, Arnheim, \emph{Nimwegen}, Cleve, Grave, \emph{Genep}, Geldern, \emph{Venloo}, Crefeld, Roermunde, \emph{Erkelens}, Jülich, \emph{Nidegen}, Münster-Eiffel, Aldenau, Mondheim, \emph{Zell}, Trarbach, Kirchberg, \emph{Oberstein}, Münchweiler, \SourceLeaf{630} Kaiserslautern, Bergzabern, Hagenau, \emph{Stollhofen}, Achern, Freydenstadt, Hornberg, Rottweil, Villingen, \emph{Geisingen}, Schaffhausen, \emph{Stein}, \emph{Frauenfeld}, Winterthur, Chur etc.

It is much to be desired that a great many observations be made near the boundary lines. They will be of value even if one simply records the observation site, whether or not a ring formed, and, if it did, how long it was visible. Since the task here is only to measure a very short interval, the lack of a seconds clock can, if necessary, be made good by counting the beats of an ordinary pocket watch or the oscillations of a heavy ball suspended from a thread, provided that in the first case the number of beats per minute is also reported and in the second case the length of the thread and the size of the ball are reported in some known unit of measure. Yet any such observation would have to be made with a good telescope if it is to be of value.

Since even at the places where the eclipse appears annular a full eighth of the solar disk remains uneclipsed, no great diminution in daylight is to be expected, much less that stars should become visible to the naked eye. For the benefit of those who might wish, during the greatest obscuration, to try to find with telescopes some of the fixed stars both nearest the Sun and brightest or to measure their distances from the lunar limb with good reflecting instruments, Prof. \textsc{Gerling} has added the following information, which may make it easier to find these stars.

\begin{center}
\begin{tabular}{lcc}
For Regulus & $P = 65\tfrac{3}{4}^\circ$ W., & $D = 17\tfrac{1}{2}^\circ$ \\
For $\ObservedBeta$ in Leo & $P = 41\tfrac{3}{4}^\circ$ E., & $D = 13^\circ$ \\
For Spica & $P = 116^\circ$ E., & $D = 36\tfrac{1}{2}^\circ$. \\
\end{tabular}
\end{center}

Here $D$ denotes the distance of the star from the center of the Sun, and $P$ the angle that the great circle from the center of the Sun to the star makes with the great circle from the center of the Sun to the north pole. The figures are calculated for $2^{\mathrm u}\ 15^{\mathrm m}$ apparent time in Paris, but they can be used throughout Germany for the time of greatest obscuration.
\par\clearpage
